\section{Introduction}\label{sec:intro}

Short-term electricity load forecasting is a foundational input to the operation
of wholesale power markets. Independent system operators rely on accurate
day-ahead demand forecasts to schedule generation through unit commitment, to
clear economic dispatch, and to size operating reserves. Because these decisions
trade off the cost of committed capacity against the risk of shortfall, even
modest reductions in forecast error translate into appreciable savings and
improved reliability across a control area. The day-ahead horizon is
particularly consequential: commitment decisions for slow-starting units are
made roughly a day in advance, so accuracy 24~hours out directly shapes the cost
and security of the resulting schedule.

We study this problem for the NEMA zone (Northeast Massachusetts/Boston) of ISO
New England, identified as load zone \texttt{.Z.NEMASSBOST} (location id 4008).
A natural and demanding reference point is the operator's own product: ISO New
England publishes an operational day-ahead reliability-region demand forecast
that market participants and operators rely upon \cite{isone}. Improving on this
forecast is a meaningful and hard bar, precisely because it is the number the
system already trusts. Our goal is therefore not merely to fit historical load
well, but to match or exceed the operational benchmark on a like-for-like basis.

In pursuing that goal we highlight and correct two methodological pitfalls that
recur in applied load-forecasting work. The first is \emph{horizon mismatch}:
comparing a model's short-lead nowcast (e.g.\ a 1-hour-ahead prediction) against
an operator's day-ahead forecast, which flatters the model by pitting an easier
task against a harder one. We instead evaluate both forecasters at the same
24-hour-ahead horizon. The second is a \emph{train/serve weather-source
mismatch}: training on weather from one provider while serving live predictions
with weather from another silently degrades deployed accuracy, because the model
learns the idiosyncrasies of the training source. We find this single mismatch
nearly doubles live day-ahead error, and that unifying the source on Open-Meteo
\cite{openmeteo} removes most of the gap to the operator.

Methodologically, we adopt a \emph{direct multi-horizon} design. Rather than
training one model and recursively rolling its predictions forward---which
compounds errors and, for load, forces a single-phase predictor to serve the
opposite phase of the diurnal cycle---we train 24 independent CatBoost
\cite{prokhorenkova2018catboost} models, one per horizon $h=1,\dots,24$, each
predicting directly from a 168-hour (one-week) lookback window. Crucially, each
horizon model receives \emph{target-hour} exogenous features: the calendar
encoding at $t{+}h$, which is known exactly, together with the forecasted
weather at $t{+}h$, the dominant driver of load a day ahead.

\paragraph{Contributions.} This paper makes the following contributions.
\begin{itemize}
  \item A direct per-horizon forecaster (\emph{Beacon}) of 24 CatBoost models
        that removes recursive error compounding, achieving on average 80.6\%
        lower MAE across the 24-hour horizon than a single-model roll-out.
  \item Target-hour exogenous features (exact calendar plus forecast weather at
        $t{+}h$) as the principal lever for day-ahead accuracy, cutting
        held-out day-ahead MAE from 183 to 77~MW.
  \item Identification and correction of a train/serve weather-source mismatch,
        which reduces live day-ahead MAE from roughly 147 to 90~MW---a
        deployment lesson rarely reported in the literature.
  \item A horizon-matched evaluation protocol against ISO New England's
        operational day-ahead forecast, paired with a rigorous leakage audit via
        a shuffled-target test.
  \item A fully reproducible, keyless, and deployable pipeline built on ISO-NE
        Web Services and Open-Meteo.
\end{itemize}

On a strictly held-out 2025 test set Beacon attains a day-ahead MAE of 76.7~MW
and an $h{=}1$ $R^2$ of 0.977, and in a horizon-matched 30-day live evaluation
it reaches a day-ahead MAE of 89.3~MW against ISO New England's 92.6~MW.

The remainder of the paper is organized as follows.
Section~\ref{sec:related} situates the work within the load-forecasting
literature. Section~\ref{sec:data} describes the load, weather, and benchmark
data. Section~\ref{sec:method} details the feature construction and the direct
multi-horizon model, and Section~\ref{sec:setup} specifies the evaluation
protocol. Section~\ref{sec:results} reports the results, Section~\ref{sec:discussion}
discusses the findings, and Sections~\ref{sec:limits} and~\ref{sec:conclusion}
cover limitations and conclusions.
